\chapter{Deviation Register}
\label{ch:deviations}

This specification is prescriptive. This appendix records every point at which
the implementation, or the published description \cite{dstm2026}, departs from
it.

Each entry carries a class:

\begin{description}[leftmargin=2.4cm,style=nextline,font=\normalfont\bfseries]
  \item[\textnormal{\textbf{soundness}}] a legal program can obtain a wrong
    answer, or an restriction that the correctness argument relies on is
    neither enforced nor documented.
  \item[\textnormal{\textbf{liveness}}] a legal program can fail to make
    progress.
  \item[\textnormal{\textbf{latent}}] a defect not reachable under current
    usage but reachable under documented usage.
  \item[\textnormal{\textbf{robustness}}] misuse produces silent corruption or
    a hang rather than a diagnostic.
  \item[\textnormal{\textbf{documentation}}] the code is correct; a comment,
    header or the paper is not.
\end{description}

\begin{table}[htbp]
\centering
\small
\begin{tabularx}{\textwidth}{@{}l l >{\raggedright\arraybackslash}X l@{}}
\toprule
\textbf{ID} & \textbf{Class} & \textbf{Summary} & \textbf{Affects} \\
\midrule
\ref{dev:zero-distance}    & soundness     & $\lag{p}{q}=0$ off-diagonal accepted but treated as absence of a channel & \cref{def:comm-graph} \\
\ref{dev:asymmetric}       & soundness     & asymmetric matrices accepted, only one row consulted & \cref{def:comm-graph} \\
\ref{dev:monotone-origin}  & soundness     & monotone origin time never enforced & \cref{res:monotone} \\
\ref{dev:write-disjoint}   & soundness     & write-disjointness never enforced & \cref{thm:determinism} \\
\ref{dev:lookahead-uncheck}& soundness     & $\delta$/lag contract unstated and unchecked & \crefrange{thm:write-closure}{thm:discard-safety} \\
\ref{dev:destroy-lifetime} & soundness     & destroy frees a recycled extent, non-collectively, while peers hold pointers & \cref{def:addressability} \\
\ref{dev:clock-width}      & latent        & clock stored as \code{int}, compared against \code{int64\_t} & \cref{def:local-clock} \\
\ref{dev:cmp-truncation}   & latent        & \code{int64\_t} difference truncated to \code{int} in comparator & --- \\
\ref{dev:darray-bounds}    & robustness    & no bounds check on darray global index & \cref{def:addressability} \\
\ref{dev:noncollective}    & robustness    & \code{darray} creation has no consensus guard & \cref{res:collective} \\
\ref{dev:non-inheriting}   & documentation & ``version history'' suggests inheritance; versions are zero-filled & \cref{def:independence} \\
\ref{dev:phase-strictness} & documentation & publish guard is $\clk{p}\ge\eff$, comment claims $>$ & \cref{res:phase} \\
\ref{dev:clock-conflation} & documentation & \code{last\_time} documented as $\max(t+\delta)$, clock tracks $t$ & \cref{res:terminate} \\
\ref{dev:barrier-doc}      & documentation & header says two-pass ring, implementation is three-round & \cref{thm:barrier-three-rounds} \\
\bottomrule
\end{tabularx}
\caption{Deviation register at a glance.}
\label{tab:deviations}
\end{table}

\clearpage

% ===========================================================================
\section{Soundness}

\begin{deviation}[Zero distance means unreachable]\label{dev:zero-distance}
\textbf{Class:} soundness.

\textbf{Statement.} (G3) of \cref{def:comm-graph} requires
$\lag{p}{q}\ge 1$ for reachable $p \ne q$. The implementation accepts
$\lag{p}{q}=0$ from the public constructor and then treats such a pair as
\emph{unreachable}: every loop that scans the \gls{neighbourhood} skips entries with
$\lag{p}{q} \le 0$.

\textbf{Evidence.} \srcref{distance.c}{655}--\srcref{distance.c}{662} validates
only $\Dmat_{ii}=0$ and $\Dmat_{ij}\ge-1$, so a zero off-diagonal passes. The
consumers then discard it: \srcref{ops.c}{42} (heartbeat broadcast),
\srcref{ops.c}{90} (publish guard), \srcref{ops.c}{154} (retire guard) all
\code{continue} on \code{dist <= 0}. The header documents the input as
``non-negative integers'' (\srcref{dstm.h}{261}) while the code also accepts
$-1$.

\textbf{Consequence.} Such a pair exchanges no \glspl{heartbeat} and imposes
no constraint on either side's promotions, so \glspl{version} publish on the
local clock alone and a remote put can land on a published one --- the abort
path of \cref{thm:write-closure} at best, a lost update at worst. The claim in
\cite{dstm2026} that ``a distance of 0 between all \glspl{rank} is a correct
setup'' is therefore false as implemented: that configuration is unsound, not
merely unoptimised. The intended reading --- all-to-all synchronisation every
timestep --- is $\lag{p}{q}=1$, which \code{dstm\_distance\_create\_complete}
already produces.

\textbf{Fix.} Reject $\lag{p}{q}=0$ for $p \ne q$ in
\code{dstm\_distance\_create\_from\_matrix}, and align the header text with
the $-1$ sentinel.

\textbf{Test.} \code{create\_from\_matrix} with a zero off-diagonal returns
\code{-DSTM\_EINVAL}.
\end{deviation}

\begin{deviation}[Asymmetric matrices are half-checked]\label{dev:asymmetric}
\textbf{Class:} soundness.

\textbf{Statement.} (G2) of \cref{def:comm-graph} requires symmetry. The
implementation neither enforces nor requires it at construction, but consults
only row \code{self} when evaluating promotion guards, so the two directions
of an asymmetric pair are never reconciled.

\textbf{Evidence.} \code{dstm\_distance\_create\_from\_matrix} performs no
symmetry check; \srcref{dstm.h}{251} explicitly states entries ``do not need
to be symmetric''. Both guards read $\lag{\text{self}}{i}$ only:
\srcref{ops.c}{89} and \srcref{ops.c}{152}. \cite{dstm2026} states
``distances are assumed symmetrical''.

\textbf{Consequence.} With $\lag{p}{q} < \lag{q}{p}$, $\proc{q}$ retires
\glspl{version} on its own larger lag while $\proc{p}$ may still read
further back, so a legal get can target an evicted version and
\cref{thm:discard-safety} fails. \code{dstm\_dgraph\_derive\_distance}
symmetrises explicitly, which is why the dgraph path has not exposed this.

\textbf{Fix.} Enforce symmetry at construction, or replace both guard lookups
with $\max(\lag{p}{q}, \lag{q}{p})$ and drop (G2).

\textbf{Test.} \code{create\_from\_matrix} with an asymmetric pair returns
\code{-DSTM\_EINVAL}.
\end{deviation}

\begin{deviation}[Monotone origin time is not enforced]\label{dev:monotone-origin}
\textbf{Class:} soundness.

\textbf{Statement.} \Cref{res:monotone} requires each process to issue \glspl{operation}
with nondecreasing \gls{origin-time}. Nothing checks it.

\textbf{Evidence.} The clock update is guarded by
\code{if (region->local\_heartbeat\_time < time)}
(\srcref{ops.c}{603}, \srcref{ops.c}{672}), so an operation with a smaller
origin time is accepted silently and simply fails to move the clock.
\cite{dstm2026} states the rule --- ``a process is only allowed to issue
operations with a monotonically increasing originating time'' --- as a
property of the system rather than as a program restriction.

\textbf{Consequence.} \Cref{thm:write-closure} argues that a process whose
clock has passed $\eff$ can no longer produce a put at $\eff$. A program going
backwards in origin time invalidates that directly, and can write to a
published \gls{version}.

\textbf{Fix.} Reject $t < \clk{p}$ with \code{-DSTM\_EDOM}, or assert under a
debug build.
\end{deviation}

\begin{deviation}[Write-disjointness is not enforced]\label{dev:write-disjoint}
\textbf{Class:} soundness.

\textbf{Statement.} \Cref{res:write-disjoint} requires that no two writes
target the same \gls{cell}. Nothing checks it, and the model has no fallback rule
for when it is violated.

\textbf{Evidence.} \cite{dstm2026} states ``as a simplification, overlapping
writes are ignored (not a conflict)''. The apply path
(\srcref{ops.c}{254}, \srcref{ops.c}{311}) is an unconditional \code{memcpy}
into the version buffer, so the surviving value is whichever write the
progress loop applies last --- a function of message arrival order.

\textbf{Consequence.} \Cref{thm:determinism} fails: the same program on the
same input can produce different results on different runs. This is the one
deviation that silently breaks reproducibility --- for a simulation framework,
the property users are least willing to lose.

\textbf{Fix.} Under a debug build, record the written byte ranges per version
and abort on overlap. The cost is proportional to the number of \glspl{put} per
version and is acceptable outside production builds. Note that the mailbox
idiom (\cref{sec:gloss-idioms}) discharges the restriction structurally, so
conforming applications pay nothing in release builds.
\end{deviation}

\begin{deviation}[The delta/lag contract is unstated]\label{dev:lookahead-uncheck}
\textbf{Class:} soundness.

\textbf{Statement.} \Cref{res:lookahead} ($\delta \ge \lag{p}{q}$ for remote
puts) and \cref{res:staleness} ($|\delta| \le \lag{p}{q}$ for remote \glspl{get})
are the assumptions that make both promotion guards sound. Neither appears in
the header documentation, in the source comments, or in \cite{dstm2026}, and
neither is checked.

\textbf{Evidence.} The only \gls{delta} validation is the sign test:
$\delta \ge 0$ for put (\srcref{ops.c}{592}), $\delta \le 0$ for get
(\srcref{ops.c}{661}), plus the local-process rule for $\delta = 0$
(\srcref{ops.c}{598}, \srcref{ops.c}{667}). The guards that depend on the
contract are \srcref{ops.c}{99} and \srcref{ops.c}{163}.

\textbf{Consequence.} Violating \cref{res:lookahead} lets a put arrive at a
published \gls{version}. Violating \cref{res:staleness} targets an evicted
one, which does not fail: the version lookup returns nothing, the request
never completes, and the caller blocks forever inside \code{dstm\_get}. A
\gls{delta} smaller than the \gls{lag} therefore hangs with no
diagnostic.

\textbf{Fix.} Document both bounds in \code{dstm.h}; check them in
\code{dstm\_put} and \code{dstm\_get}, returning \code{-DSTM\_EDOM}. The check
is two comparisons against a matrix lookup already performed nearby.

\textbf{Test.} A remote put with $0 < \delta < \lag{p}{q}$ and a remote get
with $|\delta| > \lag{p}{q}$ both return \code{-DSTM\_EDOM}.
\end{deviation}

\begin{deviation}[Destruction frees a recycled extent, non-collectively]
\label{dev:destroy-lifetime}
\textbf{Class:} soundness.

\textbf{Statement.} Destroying a derived structure releases its extent
immediately and purely locally, while every peer still holds the base offset
that was published at creation. Since freed extents are recycled, a subsequent
structure can occupy the same addresses, and a peer reading through the old
base offset then observes the new structure's bytes.

\textbf{Evidence.} \code{dstm\_darray\_destroy} calls
\code{dstm\_region\_free} on the local slice (\srcref{darray.c}{186}), as does
\code{dstm\_dgraph\_destroy} (\srcref{dgraph.c}{884}). Neither is collective,
and neither is documented as such --- contrast the creation paths, which are
collective and exchange the base offsets by all-gather
(\srcrefr{darray.c}{128}{129}). Freed extents return to the free list and are
reused (\srcref{region.c}{632}).

\textbf{Consequence.} Existing \glspl{version} are insulated, because each is
an independent snapshot (\cref{conv:no-inherit}) and retains the old
structure's bytes. The hazard needs a version written \emph{after}
reallocation, which the ordinary sequence ``destroy $A$, create $B$, keep
ticking'' produces: a peer that has not destroyed its own handle for $A$, and
reads at an effective time past $B$'s first write, resolves through $A$'s
stale base offset into $B$'s extent. The read is addressable
(\cref{def:addressability}) and returns a well-defined byte --- of the wrong
structure. Nothing detects it.

\textbf{Fix.} State the lifetime condition in \code{dstm.h}: destruction is
collective, and may only occur once the structure is quiescent, which in
practice means after \code{dstm\_barrier}. Optionally, withhold freed extents
from reuse until the next barrier, which makes the safe case cheap and the
unsafe case a deterministic failure rather than silent corruption.

\textbf{Test.} Create $A$, barrier, destroy $A$ on all processes, create $B$,
and check that a peer's stale handle for $A$ is rejected rather than resolving
into $B$.
\end{deviation}

% ===========================================================================
\section{Latent defects}

\begin{deviation}[Clock stored at the wrong width]\label{dev:clock-width}
\textbf{Class:} latent.

\textbf{Statement.} The \gls{local-clock} is declared \code{int} but holds a
\code{dstm\_time\_t}, which is \code{int64\_t}.

\textbf{Evidence.} \code{int local\_heartbeat\_time;} at
\srcref{region.h}{158}, compared against and assigned from
\code{dstm\_time\_t} at \srcref{ops.c}{603}--\srcref{ops.c}{605} and
\srcref{ops.c}{672}--\srcref{ops.c}{674}.

\textbf{Consequence.} Truncation beyond $2^{31}$: the clock appears to move
backwards, \glspl{heartbeat} stop, and neighbours stall --- a liveness failure
presenting as a hang after a long run.

\textbf{Fix.} Change the field to \code{dstm\_time\_t}.
\end{deviation}

\begin{deviation}[Comparator truncates a 64-bit difference]\label{dev:cmp-truncation}
\textbf{Class:} latent.

\textbf{Statement.} \code{dstm\_request\_cmp\_departure} computes a
\code{dstm\_time\_t} difference and returns it as \code{int}.

\textbf{Evidence.} \code{ta = a->time - b->time; if (ta) return ta;} in
\code{src/request.c}, with return type \code{int}.

\textbf{Consequence.} For differences whose low 32 bits are zero the
comparator reports equality, and otherwise can report the wrong sign,
mis-ordering the pending-request list. Ordering is a scheduling heuristic, not
a correctness requirement, which is why this is latent.

\textbf{Fix.} Return a width-safe three-way comparison, and likewise for the
\gls{delta} tiebreak:
\begin{lstlisting}
return (a->time > b->time) - (a->time < b->time);
\end{lstlisting}
\end{deviation}

% ===========================================================================
\section{Robustness}

\begin{deviation}[Unchecked darray indices and partitions]\label{dev:darray-bounds}
\textbf{Class:} robustness.

\textbf{Statement.} \code{dstm\_darray\_get} and \code{dstm\_darray\_put}
never validate \code{global\_index} against the total element count, and
\code{dstm\_darray\_create\_equal} can produce a \gls{partition} that diverges
across \glspl{process}.

\textbf{Evidence.} The owner search clamps to $[0, \worldsize-1]$
(\srcrefr{darray.c}{38}{51}) so an out-of-range index resolves to the last
process at an offset past the end of its slice, with no error. Separately, when
\code{total\_elements < world\_size} the tail processes compute
\code{local\_elements == 0} and return \code{-DSTM\_EINVAL} at
\srcref{darray.c}{73}, \emph{before} the \code{MPI\_Allgather} at
\srcref{darray.c}{120} that the other processes have already entered.

\textbf{Consequence.} The resolved offset lies outside the darray's extent, so
the operation is not addressable (\cref{def:addressability}). Nothing faults:
a \gls{region} is one contiguous \gls{address-space}, so the offset simply
lands in some other handle's extent and reads or corrupts another structure's
data. The partition case is a hang.

\textbf{Fix.} Bounds-check against \code{total\_elements}; reject
\code{total\_elements < world\_size} uniformly on all processes.
\end{deviation}

\begin{deviation}[Non-collective error paths in darray creation]\label{dev:noncollective}
\textbf{Class:} robustness.

\textbf{Statement.} \code{dstm\_darray\_create} contains early returns on
allocation failure that are not agreed across processes, so one process's local
failure desynchronises the collective.

\textbf{Evidence.} Failure returns at \srcref{darray.c}{85},
\srcref{darray.c}{90}, \srcref{darray.c}{96}, \srcref{darray.c}{102} and
\srcref{darray.c}{112} all precede or interleave with the two
\code{MPI\_Allgather} calls. Contrast \code{dstm\_dgraph\_create}, which
guards every collective with \code{dstm\_msg\_consensus}
(\srcref{dgraph.c}{537} and following) precisely to avoid this.

\textbf{Consequence.} A single out-of-memory process leaves every other process
blocked in \code{MPI\_Allgather}.

\textbf{Fix.} Apply the \code{dstm\_msg\_consensus} pattern already used by
dgraph.
\end{deviation}

% ===========================================================================
\section{Documentation}

\begin{deviation}[Versions do not inherit]\label{dev:non-inheriting}
\textbf{Class:} documentation.

\textbf{Statement.} \cite{dstm2026} describes versioning as ``each \gls{rank}
retains several \glspl{version} (i.e.\ a \gls{history}) of its state'', which suggests a
version is a snapshot of evolving state. It is not: a new version is
zero-filled and inherits nothing (\cref{def:independence}).

\textbf{Evidence.} \code{dstm\_region\_version\_find\_or\_create} allocates
with \code{calloc(1, region->size)} and does not copy any predecessor; the
accompanying comment states the zero-fill is relied upon by the
``publish only what you change'' pattern.

\textbf{Consequence.} No implementation defect. The risk is user-facing: a
program that writes a value at $\eff$ and reads the same address at $\eff+1$
without rewriting it silently reads zero. This is the single most likely
first-use error, and it is not documented anywhere in \code{dstm.h}.

\textbf{Fix.} State the semantics in the header, next to
\code{dstm\_darray\_put}. See \cref{rem:republish}.
\end{deviation}

\begin{deviation}[Publish guard is non-strict]\label{dev:phase-strictness}
\textbf{Class:} documentation.

\textbf{Statement.} The comment above the local-clock clause of the publish
guard argues that a version must stay writable ``until our own clock has moved
strictly past \code{v->time}''. The code implements $\clk{p} \ge \eff$, not
$\clk{p} > \eff$.

\textbf{Evidence.} Comment at \srcrefr{ops.c}{103}{111}; the test at
\srcref{ops.c}{112} is \code{if (region->local\_heartbeat\_time < v->time)}.

\textbf{Consequence.} None: the code is right and the comment wrong, but for a
reason worth recording. The non-strict guard leaves exactly one gap --- a
local put with $\delta = 0$ at $t = \clk{p}$, still legal after publication.
\Cref{prop:phase-minimality} shows this is the only gap and that
\cref{res:phase} closes precisely it. Strengthening the guard to $\clk{p} >
\eff$ would make \cref{res:phase} unnecessary but also make it impossible for
a process to read its own current tick --- the whole \glslink{push-model}{push
idiom}.

\textbf{Fix.} Replace the comment with a reference to \cref{res:phase}.
\end{deviation}

\begin{deviation}[Clock tracks origin, documentation says effective]
\label{dev:clock-conflation}
\textbf{Class:} documentation.

\textbf{Statement.} \code{dstm\_fini} and \code{dstm\_barrier} document their
\code{last\_time} argument as ``the highest effective timestep submitted by
this process, i.e.\ $\max(\mathit{time} + \mathit{delta})$''. The clock that the
version state machine actually consults advances to $\mathit{time}$ only.

\textbf{Evidence.} \srcref{dstm.h}{337} and \srcref{dstm.h}{359} for the
documented meaning; \srcref{ops.c}{603}--\srcref{ops.c}{605} for the update,
which uses \code{time} and not \code{time + delta}.

\textbf{Consequence.} None for a conforming program, because the \gls{barrier}'s
drained flag bypasses the guards that would care. But the two quantities are
genuinely different --- they differ by up to the largest \gls{delta} in use --- and
conflating them in the documentation makes \cref{res:terminate} hard to
discharge correctly. The restriction as stated here uses the \gls{effective-time},
matching the header rather than the clock.

\textbf{Fix.} Name the two quantities distinctly in the header.
\end{deviation}

\begin{deviation}[Barrier documented as two-pass]\label{dev:barrier-doc}
\textbf{Class:} documentation.

\textbf{Statement.} \code{dstm\_barrier} is documented as ``a two-pass ring
token algorithm''. The implementation is a three-round protocol, and
\cref{thm:barrier-three-rounds} shows three rounds are necessary.

\textbf{Evidence.} \srcref{dstm.h}{348} says two-pass; the protocol comment at
\srcrefr{barrier.c}{19}{56} and the state machine at
\srcrefr{barrier.c}{113}{225} implement a first round, a second round and an exit round.

\textbf{Fix.} Correct the header; cite the necessity argument.
\end{deviation}
